\section{Introduction}
\label{sec:intro}

Precise 3D medical image segmentation is critical for clinical applications such as disease screening, surgical planning, and image-guided intervention. To meet these demands, deep learning architectures have historically relied on the encoder-decoder paradigm established by U-Net \cite{ronneberger2015u,cicek20163d}. Over the past several years, the field has witnessed a continuous evolution toward increasingly expressive encoder representations. Concurrently, modern architectures—including UNETR \cite{hatamizadeh2022unetr}, Swin UNETR \cite{hatamizadeh2022swin}, 3D UX-Net \cite{lee20233d}, and MedNeXt \cite{roy2023mednext}—have largely retained dense, multi-resolution decoding stages. More recent designs further enhance decoder capability through advanced feature fusion, refinement, and upsampling modules. This trajectory reflects an implicit assumption that heavy decoder computation is universally beneficial for accurate segmentation. 

Despite these continuous advances, existing literature evaluates decoder design almost exclusively through aggregate, end-point metrics such as the Dice similarity coefficient or Hausdorff distance. Such global metrics offer a macroscopic view of model performance and efficiency trade-offs, but they critically obscure the micro-level behavior of the network. Specifically, they fail to answer a fundamental scientific question: \textit{Where, how much, and under what conditions is additional decoder computation actually useful?} 

To address this gap, we present a systematic characterization study of decoder computation in 3D medical image segmentation. Rather than introducing a new architecture, we employ matched pairs of full and efficient decoders (specifically building upon the EffiDec3D framework \cite{effidec3d}) as controlled experimental instruments. By isolating the decoder capacity while holding the encoder, data, and training protocols constant, we directly observe the voxel-level disagreements---the ``flips''---caused by adding decoder capacity. 

Our characterization framework yields the following core contributions that challenge the prevailing assumption of uniform decoder utility:
\begin{itemize}
    \item \textbf{Heterogeneous Contribution}: We demonstrate that the marginal benefit of a full decoder over an efficient baseline is almost entirely localized to thin boundary bands and small, difficult anatomical structures.
    \item \textbf{Concentrated Benefit}: Our oracle analysis reveals that the vast majority of the full decoder's positive corrections are clustered in a small fraction of the image volume ($<5\%$), providing direct evidence of structured decoder redundancy.
    \item \textbf{Predictability of Instability}: We show that while lightweight inference-time uncertainty signals (e.g., entropy) can reliably predict \textit{where} predictions will be unstable, they fail to predict the \textit{direction} of improvement.
    \item \textbf{Net-Neutral Voxel Correctness}: At the voxel level, the matched full decoder's net benefit is statistically indistinguishable from zero, representing a net-neutral effect where positive corrections are counterbalanced by negative degradations.
\end{itemize}

Collectively, these findings demonstrate that modern decoder computation exhibits a structured pattern of heterogeneous contribution, concentrated benefit, and partially predictable instability, providing the scientific foundation for future adaptive and dynamic decoding strategies. The rest of this paper is organized as follows: Section \ref{sec:related} reviews related work, Section \ref{sec:methods} details our characterization framework, Section \ref{sec:experiments} presents the experimental setup and results, Section \ref{sec:discussion} discusses the implications, and Section \ref{sec:conclusion} concludes the paper.

\begin{figure}[t]
\centering
\includegraphics[width=\textwidth]{figures/obs-uxnet-concat/figure1_motivation.png}
\caption{Motivation: Decoder computation yields spatially localized flips. While the full decoder corrects some efficient decoder errors (positive flips), it simultaneously degrades other previously correct predictions (negative flips) at high-entropy boundaries. This illustrates that decoder computation is not uniformly beneficial.}
\label{fig:motivation}
\end{figure}
